\documentclass{article}

\usepackage[english]{babel}
\usepackage[letterpaper,top=2cm,bottom=2cm,left=3cm,right=3cm,marginparwidth=1.75cm]{geometry}
\setlength{\parindent}{0pt}

\usepackage{amsmath}
\usepackage{amsthm}
\usepackage{amssymb}
\usepackage{graphicx}
\usepackage{float}
\usepackage{listings}
\usepackage[colorlinks=true, allcolors=blue]{hyperref}
\usepackage{xcolor}

\lstset{
  basicstyle=\ttfamily\small,
  numbers=left,
  numberstyle=\tiny,
  stepnumber=1,
  numbersep=8pt,
  backgroundcolor=\color{gray!10},
  frame=single,
  rulecolor=\color{gray!60},
  tabsize=2,
  captionpos=b,
  breaklines=true,
  breakatwhitespace=true,
  showstringspaces=false,
  keywordstyle=\color{blue},
  commentstyle=\color{green!50!black},
  stringstyle=\color{orange},
}

\DeclareMathOperator{\diag}{diag}
\DeclareMathOperator*{\argmin}{arg\,min}

\newcommand{\EE}{\mathbb{E}}
\newcommand{\PP}{\mathbb{P}}
\newcommand{\cC}{\mathcal{C}}
\newcommand{\cD}{\mathcal{D}}
\newcommand{\cW}{\mathcal{W}}
\newcommand{\one}{\mathbf{1}}

\title{DSC 291: Learning Theory Assignment 8}
\author{Tongle Shen}

\begin{document}
\maketitle

\section{Part A: Choosing Regularization by Validation}

\subsection{Problem 1. Validation Selection for a Finite Candidate Set}

Condition on fixed predictors
$$
h_1,\dots,h_K
$$
that are independent of the validation set $V$. Define
$$
\Delta_k := L_{\cD}(h_k)-L_V(h_k).
$$
Let
$$
k^\star\in \argmin_{k\in[K]} L_{\cD}(h_k),
\qquad
\widehat{k}\in \argmin_{k\in[K]} L_V(h_k).
$$
Then
$$
L_{\cD}(h_{\widehat{k}})
=
L_V(h_{\widehat{k}})+\Delta_{\widehat{k}}
\le
L_V(h_{k^\star})+\Delta_{\widehat{k}}
=
L_{\cD}(h_{k^\star})-\Delta_{k^\star}+\Delta_{\widehat{k}}.
$$
Therefore
$$
L_{\cD}(h_{\widehat{k}})
\le
\min_{k\in[K]}L_{\cD}(h_k)+2\max_{k\in[K]}|\Delta_k|.
$$
It remains to bound the expected maximum validation deviation.

For fixed $k$, $L_V(h_k)$ is the average of $n_V$ independent $[0,1]$ losses with mean $L_{\cD}(h_k)$. Hoeffding's lemma gives, for every $s>0$,
$$
\EE_V e^{s\Delta_k}\le e^{s^2/(8n_V)},
\qquad
\EE_V e^{-s\Delta_k}\le e^{s^2/(8n_V)}.
$$
Hence
$$
\EE_V e^{s|\Delta_k|}
\le
\EE_V e^{s\Delta_k}+\EE_V e^{-s\Delta_k}
\le
2e^{s^2/(8n_V)}.
$$
Using the exponential maximum bound,
$$
\EE_V\max_k |\Delta_k|
\le
\frac1s\log\left(\sum_{k=1}^K \EE_V e^{s|\Delta_k|}\right)
\le
\frac{\log(2K)}{s}+\frac{s}{8n_V}.
$$
Optimizing at
$$
s=\sqrt{8n_V\log(2K)}
$$
gives
$$
\EE_V\max_k |\Delta_k|
\le
\sqrt{\frac{\log(2K)}{2n_V}}.
$$
Taking expectation in the earlier inequality yields
$$
\EE_V L_{\cD}(h_{\widehat{k}})
\le
\min_{k\in[K]}L_{\cD}(h_k)
+
2\sqrt{\frac{\log(2K)}{2n_V}}.
$$

\subsection{Problem 2. Oracle Inequality over a Grid}

Condition on the training sample $T$. Then the trained predictors
$$
\{h_\lambda:\lambda\in\Lambda\}
$$
are fixed with respect to the validation set $V$, so Problem 1 gives
$$
\EE_V\!\left[L_{\cD}(h_{\widehat{\lambda}})\mid T\right]
\le
\min_{\lambda\in\Lambda}L_{\cD}(h_\lambda)
+
2\sqrt{\frac{\log(2K)}{2n_V}}.
$$
Now take expectation over $T$:
$$
\EE L_{\cD}(h_{\widehat{\lambda}})
\le
\min_{\lambda\in\Lambda}\EE_T L_{\cD}(A_\lambda(T))
+
2\sqrt{\frac{\log(2K)}{2n_V}}.
$$

For each fixed $\lambda$, the Week 8 stable RERM theorem gives, for every $u\in\cW$,
$$
\EE_T L_{\cD}(A_\lambda(T))
\le
L_{\cD}(u)+\lambda\Psi(u)+\frac{2G^2}{\lambda\alpha n_T}.
$$
Taking the infimum over $u$ and then the minimum over $\lambda\in\Lambda$ gives
$$
\EE L_{\cD}(h_{\widehat{\lambda}})
\le
\min_{\lambda\in\Lambda}\inf_{u\in\cW}
\left\{
L_{\cD}(u)+\lambda\Psi(u)+\frac{2G^2}{\lambda\alpha n_T}
\right\}
+
2\sqrt{\frac{\log(2K)}{2n_V}}.
$$

The validation term is only logarithmic in $K$ because the validation set is used to choose among $K$ already-trained candidates. Uniformly controlling the $K$ validation estimates costs
$$
O\!\left(\sqrt{\frac{\log K}{n_V}}\right),
$$
not $O(\sqrt{K/n_V})$.

\subsection{Problem 3. Adapting to an Unknown Comparator Scale}

Fix a comparator $u$ and write
$$
a:=\Psi(u).
$$
The $\lambda$-dependent part of the bound from Problem 2 is
$$
f(\lambda):=\lambda a+\frac{2G^2}{\lambda\alpha n_T}.
$$
Set
$$
c:=\frac{2G^2}{\alpha n_T}.
$$
Then
$$
f(\lambda)=a\lambda+\frac{c}{\lambda}.
$$
The continuous optimizer is
$$
\lambda_{\mathrm{opt}}(a)=\sqrt{\frac{c}{a}}
=
\sqrt{\frac{2G^2}{\alpha a n_T}},
$$
and the optimal value is
$$
f(\lambda_{\mathrm{opt}})
=
2\sqrt{ac}
=
2G\sqrt{\frac{2a}{\alpha n_T}}.
$$

By the grid assumption, choose $\lambda\in\Lambda$ with
$$
\frac{\lambda_{\mathrm{opt}}}{2}\le \lambda\le 2\lambda_{\mathrm{opt}}.
$$
Write $\lambda=t\lambda_{\mathrm{opt}}$ with $t\in[1/2,2]$. Since $a\lambda_{\mathrm{opt}}=c/\lambda_{\mathrm{opt}}=\sqrt{ac}$,
$$
f(\lambda)
=
\sqrt{ac}\left(t+\frac1t\right).
$$
On $[1/2,2]$,
$$
t+\frac1t\le \frac52.
$$
Therefore
$$
f(\lambda)
\le
\frac52\sqrt{ac}
=
\frac{5\sqrt{2}}{2}\,
G\sqrt{\frac{a}{\alpha n_T}}.
$$
Plugging this choice into the oracle inequality gives, for every $u$ satisfying
$$
B_{\min}^2\le \Psi(u)\le B_{\max}^2,
$$
the bound
$$
\EE L_{\cD}(h_{\widehat{\lambda}})
\le
L_{\cD}(u)
+
\frac{5\sqrt{2}}{2}\,
G\sqrt{\frac{\Psi(u)}{\alpha n_T}}
+
2\sqrt{\frac{\log(2K)}{2n_V}}.
$$

\subsection{Problem 4. A Dyadic Grid and the Price of Tuning}

For
$$
a\in[B_{\min}^2,B_{\max}^2],
$$
the optimal regularization parameter ranges over
$$
\lambda_{\mathrm{opt}}(a)
\in
\left[
\frac{\sqrt{2}G}{B_{\max}\sqrt{\alpha n_T}},
\frac{\sqrt{2}G}{B_{\min}\sqrt{\alpha n_T}}
\right].
$$
Define
$$
\lambda_0:=\frac{\sqrt{2}G}{B_{\max}\sqrt{\alpha n_T}}
$$
and use the dyadic grid
$$
\Lambda
=
\{\lambda_0 2^j: j=0,1,\dots,J\},
\qquad
J:=\left\lceil \log_2\!\left(\frac{B_{\max}}{B_{\min}}\right)\right\rceil.
$$
Then
$$
K=J+1
=
\left\lceil \log_2\!\left(\frac{B_{\max}}{B_{\min}}\right)\right\rceil+1.
$$
For every $\lambda_{\mathrm{opt}}$ in the above interval, the first grid point at least as large as $\lambda_{\mathrm{opt}}$ is at most $2\lambda_{\mathrm{opt}}$, and is certainly at least $\lambda_{\mathrm{opt}}/2$. Thus the grid satisfies the required property.

Now assume $n$ is even and set
$$
n_T=n_V=\frac n2.
$$
The result from Problem 3 becomes
$$
\EE L_{\cD}(h_{\widehat{\lambda}})
\le
L_{\cD}(u)
+
5G\sqrt{\frac{\Psi(u)}{\alpha n}}
+
2\sqrt{\frac{\log(2K)}{n}}.
$$
The first statistical term,
$$
5G\sqrt{\frac{\Psi(u)}{\alpha n}},
$$
is the RERM learning term. The second term,
$$
2\sqrt{\frac{\log(2K)}{n}},
$$
is the validation-selection term.

If the comparator scale $\Psi(u)$ were known in advance, one could tune $\lambda$ directly to that scale and avoid selecting over a grid. Not knowing the scale introduces the validation term involving
$$
\log K
=
\log\!\left(
\left\lceil \log_2(B_{\max}/B_{\min})\right\rceil+1
\right),
$$
and also a small constant-factor loss from discretizing $\lambda$.

\newpage
\section{Part B: Approximate RERM and Optimization Error}

\subsection{Problem 1. Approximate Minimizers are Close to Exact Minimizers}

For a fixed sample $S$, the objective
$$
F_S(w)=L_S(w)+\lambda\Psi(w)
$$
is $\lambda\alpha$-strongly convex because $L_S$ is convex and $\Psi$ is $\alpha$-strongly convex. Since $w_S$ is its exact minimizer,
$$
F_S(w)\ge F_S(w_S)+\frac{\lambda\alpha}{2}\|w-w_S\|^2
$$
for every $w\in\cW$.
Apply this with $w=\widetilde{w}_S$:
$$
F_S(\widetilde{w}_S)
\ge
F_S(w_S)+\frac{\lambda\alpha}{2}\|\widetilde{w}_S-w_S\|^2.
$$
The optimization guarantee gives
$$
F_S(\widetilde{w}_S)\le F_S(w_S)+\eta.
$$
Combining the two inequalities,
$$
\frac{\lambda\alpha}{2}\|\widetilde{w}_S-w_S\|^2\le \eta.
$$
Hence
$$
\|\widetilde{w}_S-w_S\|
\le
\sqrt{\frac{2\eta}{\lambda\alpha}}.
$$

\subsection{Problem 2. Stability of Approximate RERM}

Let $S$ and $S'$ differ in one example. For any test point $z$,
$$
|\ell(\widetilde{w}_S,z)-\ell(\widetilde{w}_{S'},z)|
\le
G\|\widetilde{w}_S-\widetilde{w}_{S'}\|.
$$
By the triangle inequality,
$$
\|\widetilde{w}_S-\widetilde{w}_{S'}\|
\le
\|\widetilde{w}_S-w_S\|
+
\|w_S-w_{S'}\|
+
\|w_{S'}-\widetilde{w}_{S'}\|.
$$
Problem 1 bounds the first and third terms, and the Week 8 exact movement bound gives
$$
\|w_S-w_{S'}\|
\le
\frac{2G}{\lambda\alpha n}.
$$
Therefore
$$
\|\widetilde{w}_S-\widetilde{w}_{S'}\|
\le
2\sqrt{\frac{2\eta}{\lambda\alpha}}
+
\frac{2G}{\lambda\alpha n}.
$$
Multiplying by $G$ yields
$$
|\ell(\widetilde{w}_S,z)-\ell(\widetilde{w}_{S'},z)|
\le
\frac{2G^2}{\lambda\alpha n}
+
2G\sqrt{\frac{2\eta}{\lambda\alpha}}.
$$

\subsection{Problem 3. Learning Guarantee with Optimization Error}

Let $w_S$ denote the exact RERM solution and $\widetilde{w}_S$ the approximate solution. By Lipschitzness,
$$
L_{\cD}(\widetilde{w}_S)
\le
L_{\cD}(w_S)
+
G\|\widetilde{w}_S-w_S\|.
$$
Using Problem 1,
$$
L_{\cD}(\widetilde{w}_S)
\le
L_{\cD}(w_S)
+
G\sqrt{\frac{2\eta}{\lambda\alpha}}.
$$
Taking expectation over $S$ and applying the exact stable-RERM learning theorem to $w_S$ gives, for every $u\in\cW$,
$$
\EE_S L_{\cD}(\widetilde{w}_S)
\le
L_{\cD}(u)+\lambda\Psi(u)+\frac{2G^2}{\lambda\alpha n}
+
G\sqrt{\frac{2\eta}{\lambda\alpha}}.
$$

\subsection{Problem 4. How Accurate Must the Optimizer Be?}

Assume
$$
\Psi(u)\le B^2.
$$
Use
$$
\lambda
=
\sqrt{\frac{2G^2}{\alpha B^2 n}}
=
\frac{\sqrt{2}G}{B\sqrt{\alpha n}}.
$$
Then
$$
\lambda B^2
=
\sqrt{2}\,G B\frac{1}{\sqrt{\alpha n}},
$$
and
$$
\frac{2G^2}{\lambda\alpha n}
=
\sqrt{2}\,G B\frac{1}{\sqrt{\alpha n}}.
$$
So the bound from Problem 3 becomes
$$
\EE_S L_{\cD}(\widetilde{w}_S)
\le
L_{\cD}(u)
+
2\sqrt{2}\,GB\frac{1}{\sqrt{\alpha n}}
+
G\sqrt{\frac{2\eta}{\lambda\alpha}}.
$$

We want the optimization term to satisfy
$$
G\sqrt{\frac{2\eta}{\lambda\alpha}}
\le
\frac{GB}{\sqrt{\alpha n}}.
$$
After canceling $G$ and squaring,
$$
\frac{2\eta}{\lambda\alpha}
\le
\frac{B^2}{\alpha n}.
$$
Thus it is sufficient that
$$
\eta
\le
\frac{\lambda B^2}{2n}
=
\frac{GB}{\sqrt{2\alpha}\,n^{3/2}}.
$$
Under this condition,
$$
\EE_S L_{\cD}(\widetilde{w}_S)
\le
L_{\cD}(u)
+
(2\sqrt{2}+1)\frac{GB}{\sqrt{\alpha n}}.
$$
Equivalently,
$$
\EE_S L_{\cD}(\widetilde{w}_S)
\le
\inf_{\Psi(u)\le B^2}L_{\cD}(u)
+
(2\sqrt{2}+1)\frac{GB}{\sqrt{\alpha n}}.
$$

\subsection{Problem 5. Soft-Margin Linear Classification}

For hinge loss
$$
\ell(w,(x,y))=(1-y\langle w,x\rangle)_+,
\qquad
\|x\|_2\le R,
$$
the map $w\mapsto \ell(w,(x,y))$ is $R$-Lipschitz in $\|\cdot\|_2$. Also
$$
\Psi(w)=\frac12\|w\|_2^2
$$
is $1$-strongly convex in $\|\cdot\|_2$.

To compete with all $u$ satisfying
$$
\|u\|_2\le B,
$$
we have
$$
\Psi(u)\le \frac{B^2}{2}.
$$
Using the optimized choice
$$
\lambda
=
\sqrt{\frac{2R^2}{(B^2/2)n}}
=
\frac{2R}{B\sqrt{n}},
$$
Problem 3 gives
$$
\EE_S L_{\cD}(\widetilde{w}_S)
\le
L_{\cD}(u)
+
\frac{2RB}{\sqrt{n}}
+
R\sqrt{\frac{2\eta}{\lambda}}.
$$
Substituting $\lambda=2R/(B\sqrt n)$ into the optimization term,
$$
R\sqrt{\frac{2\eta}{\lambda}}
=
R\sqrt{\frac{\eta B\sqrt n}{R}}.
$$
Thus the approximate soft-margin RERM bound is
$$
\EE_S L_{\cD}(\widetilde{w}_S)
\le
\inf_{\|u\|_2\le B}L_{\cD}(u)
+
\frac{2RB}{\sqrt{n}}
+
R\sqrt{\frac{\eta B\sqrt n}{R}}.
$$

To make the optimization-error term no larger than
$$
\frac{RB}{\sqrt n},
$$
it suffices that
$$
R\sqrt{\frac{2\eta}{\lambda}}
\le
\frac{RB}{\sqrt n}.
$$
Equivalently,
$$
\eta\le \frac{\lambda B^2}{2n}
=
\frac{RB}{n^{3/2}}.
$$
Under this condition,
$$
\EE_S L_{\cD}(\widetilde{w}_S)
\le
\inf_{\|u\|_2\le B}L_{\cD}(u)
+
\frac{3RB}{\sqrt n}.
$$

\newpage
\section{Part C: A Weighted Euclidean Geometry}

\subsection{Problem 1. Geometry Quantities}

Let
$$
Q=\diag(q_1,\dots,q_d),
\qquad
q_j>0.
$$
The $Q$-norm is
$$
\|w\|_Q=\sqrt{w^\top Qw}.
$$
Its dual norm is
$$
\|v\|_{Q,*}
=
\sqrt{v^\top Q^{-1}v}.
$$
Indeed, by Cauchy--Schwarz,
$$
\langle w,v\rangle
=
\langle Q^{1/2}w,Q^{-1/2}v\rangle
\le
\|w\|_Q\sqrt{v^\top Q^{-1}v},
$$
and equality is attained in the direction $w\propto Q^{-1}v$.

Therefore, using $|\phi_j(x)|\le r_j$,
$$
\|\phi(x)\|_{Q,*}^2
=
\sum_{j=1}^d \frac{\phi_j(x)^2}{q_j}
\le
\sum_{j=1}^d \frac{r_j^2}{q_j}.
$$

Next,
$$
\Psi_Q(w)
=
\frac12 w^\top Qw
=
\frac12\sum_{j=1}^d q_j w_j^2.
$$
Over
$$
\cC_b=\{w\in\mathbb{R}^d: |w_j|\le b_j\ \forall j\},
$$
each term is maximized by $|w_j|=b_j$. Hence
$$
\sup_{w\in\cC_b}\Psi_Q(w)
=
\frac12\sum_{j=1}^d q_j b_j^2.
$$

\subsection{\texorpdfstring{Problem 2. Fixed $Q$}{Problem 2. Fixed Q}}

The scalar loss is $g$-Lipschitz in the prediction, so for any $w,w'$,
$$
|\ell(\langle w,\phi(x)\rangle,y)-\ell(\langle w',\phi(x)\rangle,y)|
\le
g|\langle w-w',\phi(x)\rangle|.
$$
Using the $Q$ dual norm,
$$
|\langle w-w',\phi(x)\rangle|
\le
\|w-w'\|_Q \|\phi(x)\|_{Q,*}.
$$
Therefore the loss is $G_Q$-Lipschitz with respect to $\|\cdot\|_Q$, where
$$
G_Q
\le
g\sqrt{\sum_{j=1}^d \frac{r_j^2}{q_j}}.
$$
Also $\Psi_Q$ is $1$-strongly convex with respect to $\|\cdot\|_Q$, and by Problem 1,
$$
\sup_{w\in\cC_b}\Psi_Q(w)
=
\frac12\sum_{j=1}^d q_j b_j^2.
$$

Apply the Week 8 stable-RERM theorem over $\cC_b$:
$$
\EE L_{\cD}(A_\lambda(S))
\le
\inf_{w\in\cC_b} L_{\cD}(w)
+
\lambda\left(\frac12\sum_{j=1}^d q_jb_j^2\right)
+
\frac{2g^2\sum_{j=1}^d r_j^2/q_j}{\lambda n}.
$$
Let
$$
A_Q:=\sum_{j=1}^d q_jb_j^2,
\qquad
R_Q^2:=\sum_{j=1}^d \frac{r_j^2}{q_j}.
$$
The $\lambda$-dependent expression is
$$
\frac{\lambda A_Q}{2}+\frac{2g^2R_Q^2}{\lambda n}.
$$
Optimizing over $\lambda$ gives
$$
\lambda^\star=\frac{2gR_Q}{\sqrt{A_Q n}},
$$
and the optimized value is
$$
\frac{2g}{\sqrt n}\sqrt{A_QR_Q^2}.
$$
Thus
$$
\EE L_{\cD}(A_{\lambda^\star}(S))
\le
\inf_{w\in\cC_b}L_{\cD}(w)
+
\frac{2g}{\sqrt n}
\sqrt{
\left(\sum_{j=1}^d q_jb_j^2\right)
\left(\sum_{j=1}^d \frac{r_j^2}{q_j}\right)
}.
$$

\subsection{Problem 3. Best Diagonal Geometry}

We need to minimize
$$
\left(\sum_{j=1}^d q_jb_j^2\right)
\left(\sum_{j=1}^d \frac{r_j^2}{q_j}\right)
$$
over $q_j>0$. By Cauchy--Schwarz,
$$
\left(\sum_{j=1}^d b_jr_j\right)^2
=
\left(\sum_{j=1}^d (\sqrt{q_j}b_j)\left(\frac{r_j}{\sqrt{q_j}}\right)\right)^2
\le
\left(\sum_{j=1}^d q_jb_j^2\right)
\left(\sum_{j=1}^d \frac{r_j^2}{q_j}\right).
$$
Equality holds when the two vectors are proportional, i.e. when
$$
\sqrt{q_j}b_j
=
c\frac{r_j}{\sqrt{q_j}}
$$
for a constant $c>0$. Equivalently,
$$
q_j=\frac{c r_j}{b_j}.
$$
Since any $c>0$ gives the same product, such a choice is valid for all $j$.

Therefore the minimum possible product is
$$
\left(\sum_{j=1}^d b_jr_j\right)^2.
$$
Plugging this into the fixed-$Q$ bound from Problem 2, the best diagonal-geometry excess term is
$$
\frac{2g}{\sqrt n}\sum_{j=1}^d b_jr_j.
$$

\newpage
\section{Part D: AI Usage Report}

I used AI for the proof organization, algebra checking, and LaTeX polishing in this assignment. The parts where AI was most useful were the validation-selection proof in Part A, the optimization-error constants in Part B, and the Cauchy--Schwarz geometry optimization in Part C.

The suggestions I accepted were the finite-class validation proof using Hoeffding's lemma, the discretization calculation
$$
t+\frac1t\le \frac52
\qquad
\text{for }t\in[1/2,2],
$$
and the diagonal-geometry substitution
$$
q_j\propto \frac{r_j}{b_j}.
$$
These suggestions were useful because they directly matched the constants and structure requested in the assignment.

I rejected or substantially modified two suggestions. First, I did not use a high-probability validation argument, because the problem asks for an expectation bound and the cleaner expectation proof gives the exact stated form. Second, I corrected the weighted-geometry term to scale as
$$
\frac{2g}{\sqrt n}\sum_j b_jr_j,
$$
which follows from the Week 8 theorem; a version with $2g\sqrt n$ would have the wrong statistical rate.

I verified the submitted work by checking each proof against the assignment statement and the Week 8 slides. In Part A, I checked the validation term and the factor-$2$ grid discretization constants. In Part B, I checked the distinction between approximate-stability, which has a $2G\sqrt{2\eta/(\lambda\alpha)}$ term, and the expected-risk bound, which only needs one such comparison to the exact RERM output. In Part C, I checked the dual norm formula and the equality condition in Cauchy--Schwarz. Finally, I rebuilt the PDF with \verb|pdflatex| to catch LaTeX and formatting errors.

No new repo skill was necessary for this assignment. The five most recent workflow rules I relied on were:
\begin{itemize}
\item use the shared \verb|ai4lt| Conda environment before doing homework work,
\item compile with \verb|pdflatex|,
\item create a reusable statement file for each homework folder,
\item keep generated or supporting artifacts inside the assignment-local directory,
\item keep long code out of the report and reference relative paths when code is used.
\end{itemize}

\end{document}
